\section{贸易救济措施}
\begin{enumerate}
  \item 三法三条例。\zd{出口管制法}、网络安全法和数据安全法。
  \item 对外贸易法。
        \begin{enumerate}
          \item 货物、技术和服务的进出口。
          \item 不适用于港澳台单独关税区。
          \item 外贸经营者\zd{主体扩大到了个人，审批制改为登记备案制}（工商登记 \hz 其他执业手续）。
        \end{enumerate}
  \item 出口管制法。\zd{只要有足够的政治觉悟做这种题都能做得对}。
        \begin{enumerate}
          \item 出口经营者、进口商和最终用户。
          \item \zd{（军民）两用}物项。军品、\zd{核}等管制物项相关的\zd{货物、知识产权和服务}。
          \item 涉及过境转运、通运。转运，换了交通工具。通运，没有入中国的关境且没有换交通工具。
          \item 管理部门。国务院\zd{中央军委的有关部门}。
                \begin{kk}{以清单为基础的许可证制度}
                  都要申领出口许可证。
                  \begin{enumerate}
                    \item 常规出口管制清单。5 年一更。
                    \item 临时管制清单。一般有效期 2 年。
                    \item 禁止出口清单。管制物项、特地的国家和地区、特定的组织和个人。
                  \end{enumerate}
                \end{kk}
          \item 最终用户的风险管理制度。俩词，\zd{评估、核查}。
          \item 军品出口管制。军品\zd{出口专营}资格 + \zd{出口许可}证。
        \end{enumerate}
  \item \zd{网络安全与数据跨境流动}。新增。\al{滴滴}。
        \begin{enumerate}
          \item 属地原则。个人信息和重要数据原则\zd{应当在境内存储}。
          \item 保护原则。在“境外”处理数据，损害国家安全、公共利益、公民合法权益，依法追究法律责任。
          \item 数据安全审查制度。\zd{依申请 \hz 依职权双启动}机制。
        \end{enumerate}
  \item 两反一保。
  \item 反倾销。核心，必考，而且越考越细越考越难。
        \begin{enumerate}
          \item 适用条件。\zd{倾销}（出口价 \xyh 正常价）+ \zd{实质损害、威胁、产业阻碍} + 因果关系（考不着）。
                \begin{kk}{进口产品正常价值的认定有顺序要求}
                  \begin{enumerate}
                    \item 出口国/地区的可比售价。
                    \item 出口到第三国/地区的可比售价。
                    \item 原产国构成价。
                  \end{enumerate}
                  \tcblower
                  本质是出口商低价去抢夺进口国市场的行为。
                \end{kk}
          \item 调查程序。发起 \yjt 初裁 \yjt 终裁 \yjt 救济。
                \begin{enumerate}
                  \item 发起。依申请（国内产业 \dydy 总产量 \zd{50\%，半壁江山}）、依职权（商务部）。
                  \item 救济。可复议，可诉讼。联动行政法，不属于复议前置的情形。
                \end{enumerate}
          \item 措施和措施的细节。\zd{调查机关是商务部，执行机关是海关}。
                \begin{enumerate}
                  \item 初裁后终裁前。临时反倾销措施仅限临时税 \hz 提供担保。
                        \begin{enumerate}
                          \item 进口经营者承担。
                          \item 终裁确定\zd{不征}税 \hz 不追溯征税的，\zd{应退回，担保应解除}。
                          \item 确定征收的。\zd{多退少不补}。国家可以吃亏，个人不可以。
                        \end{enumerate}
                  \item 终局裁定。征收\zd{反倾销税 \hz 价格承诺} 2 选 1。
                        \begin{kk}{反倾销税}
                          \begin{enumerate}
                            \item \zd{进口经营者}缴纳。
                            \item 可适用不同的税额/率（\zd{看人下菜}），但\zd{不应超过倾销幅度}。
                            \item 一般不超过 \zd{5} 年。期满还要征税，必须重新调查。
                            \item \zd{原则上不追溯}征收，特殊情况可以（\zd{不是应当}）。
                            \item 追溯时间有要求。\yi 实施临时反倾销措施之日前 90 日。\er 采取临时反倾销措施期间。
                          \end{enumerate}
                        \end{kk}

                        \begin{kk}{价格承诺（只适用于两反）}
                          \begin{enumerate}
                            \item 出口经营者可以作出承诺。
                            \item 商务部可以接受，拒绝要说明理由。
                            \item \zd{初裁前商务部不能}主动寻求 \hz 接受价格承诺（要先计算倾销幅度）。
                          \end{enumerate}
                        \end{kk}
                \end{enumerate}
        \end{enumerate}
  \item 反补贴。产量占国内同类产品总产量 \zd{25\%}。
        \begin{enumerate}
          \item 适用条件。进口产品存在\zd{专向补贴}。出口国政府直接、间接提供\zd{财政支持 + 特定者获益}。\al{修高速公路 \cuo}。
          \item 价格承诺的主体\zd{多了出口国政府}。
          \item 其他都同反倾销。
        \end{enumerate}
  \item 保障措施。背景。应对 WTO 关税减让义务。
        \begin{enumerate}
          \item 已经有竞争性产业，不能为了建立新产业启动保障措施。
          \item 适用条件。进口产品\zd{数量增加}（不要求是否公平）+ 产业\zd{严重}损害 \hz 严重损害威胁 + \zd{因果关系}。
          \item 针对进口产品实施，不区分来源国/地区。
          \item \zd{初裁}只能采用\zd{提高关税}的方式。\zd{终裁}可以采用\zd{提高关税 \hz 数量限制}的方式。
          \item 一般不超过 \zd{4} 年，\zd{最长 10 年}。\zd{超过 1 年}的，应当在实施期间内，按固定时间间隔\zd{逐步放宽}。
          \item 初裁并非必经程序、没有司法审查程序、不可以提起行政诉讼（\zd{无救济}）。
        \end{enumerate}
        % \item 累进评价。
\end{enumerate}
